\chapter{Evaluation}\label{ch:evaluation}
%% ---------------------------------------------------------------------
%% PURPOSE: provide the evidence for each contribution claimed in the
%% introduction. An examiner will check every claim against this chapter.
%%
%% Structure:
%%   1. Research questions / hypotheses - restate them, numbered (RQ1..).
%%   2. Setup - datasets, baselines, metrics, parameters, hardware. Enough to
%%      reproduce. Justify each metric: what does a high value mean?
%%   3. Results - one section per RQ. Each starts with the answer, then the
%%      evidence (table/figure), then the interpretation.
%%   4. Discussion & threats to validity - what could make the results wrong
%%      (internal, external, construct validity). Name the limits honestly.
%%
%% Rules of thumb:
%%   - Report numbers once, in one place (a table), and refer to it. If a
%%     number is repeated in the text, abstract or slides it MUST match.
%%   - Say what the result means, not only what it is.
%%   - Match claim strength to evidence: measured > analytic > anecdotal.
%%   - Report variance / multiple runs when results are stochastic.
%% ---------------------------------------------------------------------

\section{Research Questions}\label{sec:rqs}
\begin{description}
  \item[RQ1] \todo{...}
  \item[RQ2] \todo{...}
\end{description}

\section{Experimental Setup}\label{sec:setup}
\todo{Datasets, baselines, metrics, configuration.}

\section{Results}\label{sec:results}
\begin{table}[htbp]
  \centering
  \caption{Main results. \todo{Metric definition; best value in bold.}}
  \label{tab:main-results}
  \begin{tabular}{lrr}
    \toprule
    Method & Metric A $\uparrow$ & Metric B $\downarrow$ \\
    \midrule
    Baseline       & 0.00 & 0.00 \\
    \textbf{Ours}  & \textbf{0.00} & \textbf{0.00} \\
    \bottomrule
  \end{tabular}
\end{table}
\todo{Answer RQ1, citing \autoref{tab:main-results}.}

\section{Threats to Validity}\label{sec:threats}
\todo{Internal, external and construct validity.}
